\documentclass[letterpaper,11pt]{article}

% --- PAQUETES DE IDIOMA ---
\usepackage[spanish,mexico,es-noshorthands]{babel}
\usepackage[utf8]{inputenc}
\usepackage{csquotes}
\usepackage{iftex}

\usepackage{pdfpages}

% --- PAQUETES MATEMÁTICOS ---
\usepackage{amsmath}
\usepackage{amssymb}
\usepackage{cancel}
\usepackage{mathrsfs}

% --- PAQUETES GRÁFICOS Y DE DISEÑO ---
\usepackage{graphicx}
\usepackage{float}
\usepackage{placeins}
\usepackage{subcaption}
\usepackage[table]{xcolor}
\usepackage{makecell}
\usepackage{multirow}
\usepackage{array}
\usepackage{tikz}
\usetikzlibrary{shapes.geometric, shapes.multipart, positioning, arrows.meta, calc, shadows, decorations.pathreplacing, babel, fit, backgrounds}

% --- LISTAS ---
\usepackage{enumitem}

% --- PAQUETES DE TABLAS ---
\usepackage{booktabs}
\usepackage{longtable}

% --- COLORES ---
\definecolor{xlHeader}{HTML}{EFEFEF}
\definecolor{xlSelected}{HTML}{B7B7B7}
\definecolor{xlSummary}{HTML}{CFE2F3}
\definecolor{xlComment}{HTML}{C9DAF8}
\definecolor{armblue}{RGB}{0, 113, 188}

\definecolor{csfGV}{HTML}{FFF2CC}
\definecolor{csfID}{HTML}{D9E1F2}
\definecolor{csfPR}{HTML}{E2EFDA}
\definecolor{csfDE}{HTML}{FCE4D6}
\definecolor{csfRS}{HTML}{F8CBAD}
\definecolor{csfRC}{HTML}{D0CECE}

% --- RUTAS DE IMÁGENES ---
\graphicspath{{img/}{portada_img/}}

% --- ÍNDICE CON PUNTOS ---
\usepackage{tocloft}
\renewcommand{\cftsecleader}{\cftdotfill{\cftdotsep}}

% --- CONFIGURACIÓN DE PÁGINA (GEOMETRÍA) ---
\usepackage[
    left=18mm, 
    right=18mm,
    top=30mm,
    bottom=24mm,
    headheight=77pt, 
    papersize={21.59cm,27.94cm}
]{geometry}

% --- FORMATO DE TEXTO ---
\usepackage{setspace}
\onehalfspacing
\usepackage{parskip}

% --- ENCABEZADOS Y PIES DE PÁGINA ---
\usepackage{fancyhdr}
\pagestyle{fancy}
\fancyhf{} 
\fancyhead[L]{\includegraphics[height=1.5cm]{portada_img/der.png}}
\fancyhead[R]{\includegraphics[height=1.5cm]{portada_img/izq.png}}
\fancyfoot[C]{\thepage} 
\renewcommand{\headrulewidth}{0.4pt} 
\renewcommand{\footrulewidth}{0pt}  

% --- BIBLIOGRAFÍA ---
\usepackage[backend=biber, style=apa, sortcites, url=true]{biblatex}
\addbibresource{referencias.bib}

% --- HIPERVÍNCULOS ---
\usepackage{hyperref}
\hypersetup{
    colorlinks=true,
    linkcolor=blue,
    citecolor=blue,
    filecolor=magenta,
    urlcolor=cyan
}

% Tipos de columna personalizados
\newcolumntype{L}[1]{>{\raggedright\arraybackslash}p{#1}}
\newcolumntype{C}[1]{>{\centering\arraybackslash}p{#1}}
\newcolumntype{R}[1]{>{\raggedleft\arraybackslash}p{#1}}
\newcolumntype{M}[1]{>{\centering\arraybackslash}m{#1}}

\begin{document}

\thispagestyle{empty} % Sin encabezado ni pie en la portada

% --- ENCABEZADO DE PORTADA (TABLA DE 3 COLUMNAS) ---
\begin{center}
    % La tabla usa 'm' para centrar verticalmente las imágenes y el texto
    \begin{tabular}{m{0.15\textwidth} m{0.60\textwidth} m{0.15\textwidth}}
        % Logo Izquierdo (UNAM)
        \includegraphics[width=2cm]{portada_img/der.png} & 
        
        % Título Central
        \centering\Huge\textbf{Universidad Nacional Autónoma de México} &
        
        % Logo Derecho (Facultad)
        \raggedleft\includegraphics[width=2cm]{portada_img/izq.png}
    \end{tabular}
\end{center}

\vspace{0.3cm}

% --- INFORMACIÓN DEL REPORTE ---
\begin{center}
    {\LARGE \textbf{Facultad de Ingeniería}}\\[0.8cm]

    {\LARGE \textbf{Alumno:}}\\[0.3cm]
    {\Large Lara Hernandez Angel Husiel - 320060829} \\[0.6cm]
    {\Large Lugo Manzano Rodrigo - 320206968} \\[0.6cm]
    {\Large Flores Colin Victor Jaziel - 320266083} \\[0.6cm]
    {\Large Garcia Cortes Adolfo de Jesus - 320317264} \\[0.8cm]

    {\LARGE \textbf{Seguridad informatica basica}} \\[0.8cm]

    {\LARGE Semestre: 2027-1} \\[0.8cm]
    
    {\LARGE \textbf{Proyecto 1}}\\[0.4cm]

    {\LARGE {Gestion de riesgos de ciberseguridad de VELOX Automotriz S.A. de C.V.}}\\[0.8cm]

    {\Large \textbf{Fecha de Entrega:}}\\[0.2cm]
    {\Large 14 de Noviembre del 2026}
\end{center}


\newpage
{\hypersetup{linkcolor=black}\tableofcontents}

\newpage
\setcounter{page}{1}
\singlespacing

% =========================================================
\section{Definición Empresarial}
% =========================================================

\textbf{Aetheria FinTech \& AI Solutions S.A. de C.V.} es una empresa mexicana dedicada al desarrollo de herramientas de inteligencia artificial para el sector financiero. Su servicio principal consiste en un asistente automatizado y un motor de evaluación crediticia que ayuda a financieras y cajas de ahorro a responder dudas de clientes y calificar solicitudes de préstamo en pocos segundos.

En la Tabla~\ref{tab:definicion-empresarial} se describen los datos generales de la organización, su giro, sus objetivos y los servicios que ofrece.

\vspace{0.2cm}
\begingroup
\footnotesize
\setlength{\tabcolsep}{5pt}
\renewcommand{\arraystretch}{1.32}
\noindent
\begin{longtable}{|C{1.2cm}|L{3.6cm}|L{12.2cm}|}
\caption{Datos generales de \textit{Aetheria FinTech \& AI Solutions S.A. de C.V.}}
\label{tab:definicion-empresarial} \\
\hline
\rowcolor{xlSummary}
\textbf{Inciso} & \textbf{Elemento Empresarial} & \textbf{Descripción} \\
\hline
\endfirsthead
\hline
\rowcolor{xlSummary}
\textbf{Inciso} & \textbf{Elemento Empresarial} & \textbf{Descripción} \\
\hline
\endhead
\hline
\endlastfoot

\textbf{a.} & \textbf{Sector} & Sector financiero tecnológico (FinTech) y de tecnologías de la información, enfocado en servicios automatizados de evaluación crediticia y atención a clientes mediante inteligencia artificial. \\
\hline

\textbf{b.} & \textbf{Nombre} & \textbf{Aetheria FinTech \& AI Solutions S.A. de C.V.} (nombre comercial: \textit{Aetheria AI}). \\
\hline

\textbf{c.} & \textbf{Objetivo de negocio} & Ofrecer a instituciones microfinancieras, cooperativas de ahorro y comercios en línea una plataforma para atender clientes y evaluar solicitudes de crédito de forma automática mediante el archivo de Python que es el bot, reduciendo los tiempos de respuesta y manteniendo una disponibilidad del $99.9\%$. \\
\hline

\textbf{d.} & \textbf{Activo principal} & El \textbf{modelo de inteligencia artificial y el archivo de Python que es el bot en producción}, junto con su \textbf{dataset de entrenamiento} y la \textbf{base de datos de clientes}. En estos componentes se concentran las reglas de calificación crediticia, los pesos entrenados del modelo y los datos personales y financieros de los usuarios. \\
\hline

\textbf{e.} & \textbf{Misión y Visión} & \textbf{Misión:} Brindar soluciones de inteligencia artificial confiables y seguras que agilicen el otorgamiento de créditos y la atención al cliente, cuidando en todo momento la privacidad y la integridad de la información financiera.\newline
\textbf{Visión:} Ser una empresa referente en México en el desarrollo de sistemas automatizados para servicios financieros, destacando por la estabilidad de su plataforma y el manejo seguro de los datos de sus clientes. \\
\hline

\textbf{f.} & \textbf{Modelo de negocio} & Esquema de software como servicio (SaaS) dividido en tres modalidades:\newline
1. \textit{Renta mensual empresarial:} Cobro fijo mensual a financieras y comercios por usar el asistente conversacional y el motor de calificación.\newline
2. \textit{Cobro por consulta vía API:} Tarifa por volumen de evaluaciones de crédito y peticiones procesadas por el archivo de Python que es el bot.\newline
3. \textit{Portal web para usuarios:} Sitio web alojado en el servidor frontal webserver donde los clientes consultan su estado de crédito y reciben atención directa. \\
\hline

\textbf{g.} & \textbf{Catálogo de productos y servicios} & \textbf{1. AetheriaBot (el archivo de Python que es el bot):} Asistente virtual programado en Python y alojado en el servidor interno codeserver para responder dudas de clientes, simular préstamos y preaprobar créditos.\newline
\textbf{2. Aetheria ScoreEngine:} Módulo que consulta la base de datos del modelo de IA para estimar el riesgo de impago de cada solicitante.\newline
\textbf{3. Aetheria Web Portal:} Aplicación web pública montada sobre Apache 2.4.10 en el servidor webserver que sirve como interfaz para los clientes.\newline
\textbf{4. Aetheria Custom-Train:} Servicio de ajuste y entrenamiento del modelo utilizando el dataset de la empresa sobre la infraestructura física de servidores GPU. \\
\hline
\end{longtable}
\endgroup

% =========================================================
\section{Identificación de Activos de Tecnología e Información}
% =========================================================

El inventario tecnológico de la empresa abarca tanto la infraestructura que da soporte al modelo de inteligencia artificial (la base de datos del modelo, el servidor que la aloja, el modelo mismo junto con el archivo de Python que es el bot, el dataset de entrenamiento y los servidores físicos de cómputo GPU) como los equipos de red y servidores de aplicación que permiten dar servicio hacia internet (el servidor frontal Linux webserver y el servidor interno Windows codeserver).

% ---------------------------------------------------------
\subsection{Diagrama de Red de la Organización}
% ---------------------------------------------------------

La red de la empresa está dividida en cuatro segmentos (Figura~\ref{fig:red-aetheria}):
\begin{enumerate}[noitemsep,topsep=2pt]
    \item \textbf{Zona Externa y Perímetro:} Punto de entrada para los clientes desde internet mediante conexiones HTTPS a través de Cloudflare y el firewall perimetral.
    \item \textbf{Zona Desmilitarizada (DMZ, subred 192.168.177.0/24):} Contiene el servidor Linux frontal \textbf{webserver} con IP 192.168.177.129 (interfaz eth0), donde corren \textbf{Apache 2.4.10} en el puerto 80/tcp y \textbf{OpenSSH 6.7p1} en el puerto 13337/tcp. Este equipo tiene una segunda tarjeta de red (eth1 con IP 10.0.0.78) conectada hacia la red interna.
    \item \textbf{Red Interna de Producción (subred 10.0.0.0/24):} Segmento privado donde se ubica el servidor \textbf{codeserver} (10.0.0.201) con \textbf{Windows Server 2012 R2 Standard} (puertos RDP 3389/tcp y WinRM 5985/tcp). En este servidor se ejecuta en producción \textbf{el archivo de Python que es el bot}.
    \item \textbf{Zona de Datos y Entrenamiento IA (rango 10.0.0.210 a 10.0.0.230):} Red interna donde se encuentran el \textbf{Servidor de Base de Datos} (10.0.0.210), la \textbf{Base de Datos del Modelo de IA}, el \textbf{Dataset de Entrenamiento} (10.0.0.220) y la \textbf{Infraestructura Física GPU} (10.0.0.230).
\end{enumerate}

\begin{figure}[H]
\centering
\resizebox{0.99\textwidth}{!}{%
\begin{tikzpicture}[
    font=\sffamily\scriptsize,
    zone/.style={draw=armblue, dashed, thick, rounded corners=6pt, inner sep=11pt},
    srv/.style={draw=black!80, fill=csfGV, thick, rounded corners=3pt, text width=4.6cm, minimum height=1.7cm, align=center, inner sep=4pt},
    crit/.style={draw=black!85, fill=csfRS!45, thick, rounded corners=3pt, text width=4.8cm, minimum height=1.7cm, align=center, inner sep=4pt},
    data/.style={draw=black!85, fill=csfID, thick, rounded corners=3pt, text width=4.6cm, minimum height=1.7cm, align=center, inner sep=4pt},
    ext/.style={draw=armblue, fill=xlSummary, thick, rounded corners=4pt, text width=3.3cm, minimum height=1.5cm, align=center, inner sep=4pt},
    arr/.style={-Stealth, thick, draw=black!80},
    piv/.style={-Stealth, very thick, dashed, draw=red!75!black}
]

% ==================== FILA 1: EXTERNO Y DMZ ====================
\node[ext] (internet) at (0, 0) {
    \textbf{Usuarios y Clientes}\\
    Internet Público\\
    (HTTPS 443/tcp)
};

\node[ext] (waf) at (5.0, 0) {
    \textbf{[ACT-10] WAF y Proxy CDN}\\
    Cloudflare + Firewall\\
    (TLS 1.2/1.3 $\rightarrow$ HTTP)
};

\node[srv] (webserver) at (11.2, 0) {
    \textbf{[ACT-06] webserver (Linux)}\\
    Debian GNU/Linux 8.11 (Kernel 3.16)\\
    eth0: 192.168.177.129\\
    eth1: 10.0.0.78 | SSH: 13337/tcp
};

\node[srv] (apache) at (17.4, 0) {
    \textbf{[ACT-07] Servicio Web Apache}\\
    Apache/2.4.10 (Debian) Puerto 80/tcp\\
    \textbf{[ACT-09] Respaldos Locales:}\\
    Archivo comprimido en /var/backups/
};

% ==================== FILA 2: RED INTERNA PRODUCCIÓN ====================
\node[crit] (codeserver) at (11.2, -3.9) {
    \textbf{[ACT-08] codeserver (Windows)}\\
    Windows Server 2012 R2 Standard\\
    IP Interna: 10.0.0.201\\
    RDP 3389/tcp | WinRM 5985/tcp\\
    Cuenta: developer2024
};

\node[crit] (botmodel) at (17.4, -3.9) {
    \textbf{[ACT-03] Modelo de IA}\\
    Motor \textit{Aetheria-Core LLM}\\
    Archivo de Python que es el bot\\
    Pesos neuronales y lógica crediticia
};

% ==================== FILA 3: ZONA DE DATOS Y ENTRENAMIENTO IA ====================
\node[data] (dbsrv) at (5.0, -7.8) {
    \textbf{[ACT-02] Servidor de Base de Datos}\\
    PostgreSQL / VectorDB (10.0.0.210:5432)\\
    \textbf{[ACT-01] Base de Datos con el}\\
    \textbf{Modelo de IA y Registros PII}
};

\node[data] (dataset) at (11.2, -7.8) {
    \textbf{[ACT-04] Dataset de Entrenamiento}\\
    Repositorio NAS IA (10.0.0.220)\\
    Corpus financiero y crediticio (2.4 TB)\\
    Datos de calibración del bot
};

\node[data] (gpucluster) at (17.4, -7.8) {
    \textbf{[ACT-05] Infraestructura Física IA}\\
    Clúster de Servidores GPU (10.0.0.230)\\
    Aceleradores NVIDIA A100\\
    Procesamiento y entrenamiento del bot
};

% ==================== FONDOS DE ZONAS ====================
\begin{scope}[on background layer]
    \node[zone, fill=csfGV!25, fit=(webserver) (apache), label={[anchor=south west, font=\sffamily\scriptsize\bfseries]north west:Zona Desmilitarizada: DMZ Perimetral (192.168.177.0/24)}] (dmzbox) {};
    \node[zone, fill=csfRS!15, fit=(codeserver) (botmodel), label={[anchor=south west, font=\sffamily\scriptsize\bfseries]north west:Red Interna de Producción Aislada (10.0.0.0/24)}] (intbox) {};
    \node[zone, fill=csfID!25, fit=(dbsrv) (dataset) (gpucluster), label={[anchor=south west, font=\sffamily\scriptsize\bfseries]north west:Zona de Datos y Cómputo de Entrenamiento IA (10.0.0.210 a 10.0.0.230)}] (databox) {};
\end{scope}

% ==================== CONEXIONES ====================
\draw[arr] (internet.east) -- node[above=2pt]{\tiny HTTPS 443} (waf.west);
\draw[arr] (waf.east) -- node[above=2pt]{\tiny HTTP 80} node[below=2pt]{\tiny SSH 13337} (webserver.west);
\draw[arr] (webserver.east) -- node[above=2pt]{\tiny Local} (apache.west);

% Pivoteo SSH entre webserver y codeserver
\draw[piv] (webserver.south) -- node[right=4pt, pos=0.42, text=red!75!black]{\tiny \textbf{Túnel SSH / Pivoteo (10.0.0.78 $\rightarrow$ 10.0.0.201: RDP 3389 / WinRM 5985)}} (codeserver.north);

% Relación entre codeserver y el archivo de Python que es el bot
\draw[arr] (codeserver.east) -- node[above=2pt]{\tiny Aloja} (botmodel.west);

% Conexiones hacia la Zona de Datos e IA
\draw[arr] (codeserver.west) -| node[above, pos=0.25]{\tiny Consultas SQL / Vectores (5432/tcp)} (dbsrv.north);
\draw[arr] (botmodel.south) -- node[right=4pt, pos=0.38]{\tiny Inferencia / Pesos} (gpucluster.north);
\draw[arr] (dbsrv.east) -- node[above=2pt]{\tiny Ingesta} (dataset.west);
\draw[arr] (dataset.east) -- node[above=2pt]{\tiny Entrenamiento} (gpucluster.west);

\end{tikzpicture}%
}
\caption{Diagrama de red de \textit{Aetheria FinTech \& AI Solutions S.A. de C.V.}}
\label{fig:red-aetheria}
\end{figure}

% ---------------------------------------------------------
\subsection{Registro de Activos de Información y Tecnología}
% ---------------------------------------------------------

En la Tabla~\ref{tab:registro-activos} se listan los activos de la empresa junto con su ubicación, características técnicas y las propiedades de seguridad que se deben resguardar (Confidencialidad, Integridad y Disponibilidad).

\vspace{0.2cm}
\begingroup
\scriptsize
\setlength{\tabcolsep}{4pt}
\renewcommand{\arraystretch}{1.28}
\noindent
\begin{longtable}{|C{1.2cm}|L{2.9cm}|L{2.0cm}|L{2.3cm}|L{5.05cm}|C{1.4cm}|C{1.45cm}|}
\caption{Registro de activos de información y tecnología de \textit{Aetheria FinTech \& AI Solutions}.}
\label{tab:registro-activos} \\
\hline
\rowcolor{xlSummary}
\textbf{ID} & \textbf{Nombre del Activo} & \textbf{Categoría} & \textbf{Ubicación / IP} & \textbf{Características Técnicas} & \textbf{Propiedad CIA} & \textbf{Criticidad} \\
\hline
\endfirsthead
\hline
\rowcolor{xlSummary}
\textbf{ID} & \textbf{Nombre del Activo} & \textbf{Categoría} & \textbf{Ubicación / IP} & \textbf{Características Técnicas} & \textbf{Propiedad CIA} & \textbf{Criticidad} \\
\hline
\endhead
\hline
\endlastfoot

\textbf{ACT-01} & \textbf{Base de datos con el modelo de IA y registros de clientes} & Información / Base de Datos & Zona de Datos (10.0.0.210, puerto 5432/tcp) & Guarda los vectores e hiperparámetros del modelo de IA, además del historial de solicitudes, expedientes y datos personales (PII) de 450,000 clientes. & \textbf{C, I, D} & \cellcolor{csfRS}\textbf{Crítico (5/5)} \\
\hline

\textbf{ACT-02} & \textbf{Servidor donde vive la base de datos (dbserver)} & Infraestructura / Servidor & Zona de Datos (10.0.0.210/24) & Servidor Linux con PostgreSQL y extensión vectorial pgvector que atiende las consultas de ACT-01 solicitadas desde codeserver. & \textbf{D, I, C} & \cellcolor{csfRS}\textbf{Crítico (5/5)} \\
\hline

\textbf{ACT-03} & \textbf{El modelo de Inteligencia Artificial y el archivo de Python que es el bot} & Software / Propiedad Intelectual & Servidor codeserver (10.0.0.201) & Motor de inferencia \textit{Aetheria-Core LLM} y \textbf{el archivo de Python que es el bot} en producción. Recibe las consultas de los usuarios y calcula los dictámenes de crédito. & \textbf{I, C, D} & \cellcolor{csfRS}\textbf{Crítico (5/5)} \\
\hline

\textbf{ACT-04} & \textbf{El dataset para el entrenamiento del bot} & Información / Corpus de Datos & Almacenamiento NAS IA (10.0.0.220) & Conjunto de datos de 2.4 TB con historiales crediticios, contratos y conversaciones previas que se usan para entrenar y calibrar al bot. & \textbf{C, I} & \cellcolor{csfRS}\textbf{Crítico (5/5)} \\
\hline

\textbf{ACT-05} & \textbf{Infraestructura física para procesar y entrenar al bot} & Hardware / Clúster GPU & Centro de Datos (10.0.0.230) & Servidores físicos de alto rendimiento con tarjetas GPU NVIDIA A100 y almacenamiento NVMe RAID dedicados al procesamiento y entrenamiento del modelo. & \textbf{D, I} & \cellcolor{csfDE}\textbf{Alto (4/5)} \\
\hline

\textbf{ACT-06} & \textbf{Servidor Linux Frontal (webserver)} & Infraestructura / Servidor DMZ & DMZ 192.168.177.129 (eth0) e Interna 10.0.0.78 (eth1) & Servidor con \textbf{Debian GNU/Linux 8.11 (jessie)}, kernel 3.16.0-6-686-pae (i686), cuentas root y webadmin, servicio OpenSSH 6.7p1 en el puerto 13337/tcp y rpcbind en 111/tcp. & \textbf{I, C, D} & \cellcolor{csfRS}\textbf{Crítico (5/5)} \\
\hline

\textbf{ACT-07} & \textbf{Servicio Web Apache (Apache/2.4.10)} & Servicio / Aplicación Web & Sobre webserver\newline (192.168.177.129, puerto 80/tcp) & Servidor HTTP \textbf{Apache/2.4.10 (Debian)} en el puerto 80/tcp que publica el portal web de la empresa y el formulario de conexión con el bot. & \textbf{D, I} & \cellcolor{csfDE}\textbf{Alto (4/5)} \\
\hline

\textbf{ACT-08} & \textbf{Servidor Interno Windows (codeserver)} & Infraestructura / Servidor Producción & Red Interna aislada (10.0.0.201/24) & Servidor con \textbf{Microsoft Windows Server 2012 R2 Standard} (versión 6.3.9600), usuario developer2024, servicios \textbf{RDP (3389/tcp)}, \textbf{WinRM (5985/tcp)} e IIS local. Aloja el archivo de Python que es el bot. & \textbf{C, I, D} & \cellcolor{csfRS}\textbf{Crítico (5/5)} \\
\hline

\textbf{ACT-09} & \textbf{Archivos Comprimidos de Respaldo y Credenciales} & Información / Respaldos y Secretos & webserver (/root/ y /var/backups/) & Copias de respaldo comprimidas del proyecto guardadas localmente en webserver, las cuales contienen scripts de configuración y el archivo de texto con las contraseñas. & \textbf{C, I} & \cellcolor{csfDE}\textbf{Alto (4/5)} \\
\hline

\textbf{ACT-10} & \textbf{Perímetro WAF/CDN, Certificados TLS y DNS Corporativo} & Seguridad / Red Perimetral & Perímetro Externo (Cloudflare / Firewall) & Configuración del dominio, certificados digitales X.509 (SHA-256 con RSA), protocolos TLS y filtrado de tráfico entrante. & \textbf{D, I, A} & \cellcolor{csfDE}\textbf{Alto (4/5)} \\
\hline
\end{longtable}
\endgroup

% =========================================================
\section{Registro de Amenazas Relevantes}
% =========================================================

El análisis de amenazas de la empresa contempla tanto las vulnerabilidades técnicas detectadas en los servidores y componentes de red como los perfiles de atacantes internos y externos que tienen interés en comprometer el sistema.

% ---------------------------------------------------------
\subsection{Vulnerabilidades Técnicas en los Activos}
% ---------------------------------------------------------

\begingroup
\scriptsize
\setlength{\tabcolsep}{4pt}
\renewcommand{\arraystretch}{1.28}
\noindent
\begin{longtable}{|C{2.3cm}|C{1.4cm}|L{3.6cm}|L{7.2cm}|C{1.8cm}|}
\caption{Vulnerabilidades técnicas presentes en los activos de \textit{Aetheria FinTech \& AI Solutions}.}
\label{tab:vulnerabilidades} \\
\hline
\rowcolor{xlSummary}
\textbf{CVE / CWE} & \textbf{Activo} & \textbf{Vulnerabilidad / Debilidad} & \textbf{Detalle Técnico en la Infraestructura} & \textbf{CVSS v3.1} \\
\hline
\endfirsthead
\hline
\rowcolor{xlSummary}
\textbf{CVE / CWE} & \textbf{Activo} & \textbf{Vulnerabilidad / Debilidad} & \textbf{Detalle Técnico en la Infraestructura} & \textbf{CVSS v3.1} \\
\hline
\endhead
\hline
\endlastfoot

\textbf{CVE-2021-4034} (\textit{PwnKit}) / CWE-787 & \textbf{ACT-06} (webserver) & Escalación local de privilegios en pkexec (Polkit) por escritura fuera de límites. & El binario SUID /usr/bin/pkexec instalado en Debian 8.11 no valida correctamente el número de argumentos y permite al usuario webadmin obtener acceso como root (uid=0) de inmediato. & \cellcolor{csfDE}\textbf{7.8 (Alto)} \\
\hline

\textbf{CVE-2016-5195} (\textit{DirtyCow}) / CWE-362 & \textbf{ACT-06} (webserver) & Condición de carrera (\textit{Race Condition}) en el mecanismo \textit{Copy-on-Write} del kernel Linux 3.16.0-6. & La versión del kernel 3.16.0-6-686-pae sin actualizar permite modificar páginas de memoria de solo lectura para escalar privilegios a root. & \cellcolor{csfDE}\textbf{7.8 (Alto)} \\
\hline

\textbf{CVE-2017-0358} / CWE-269 & \textbf{ACT-06} (webserver) & Carga insegura de módulos en ntfs-3g con bit SUID activo. & El paquete ntfs-3g instalado en Debian 8 conserva el permiso SUID activo y representa otra vía de escalación local de privilegios. & \cellcolor{csfDE}\textbf{7.8 (Alto)} \\
\hline

\textbf{CWE-521} / \textbf{CWE-307} & \textbf{ACT-06} (webserver) & Contraseñas débiles y falta de bloqueo por intentos fallidos en SSH (13337/tcp). & La cuenta webadmin utiliza una contraseña de diccionario (kitty123) vulnerable a ataques con Hydra, mientras que el hash SHA-512 de root corresponde a una clave débil (colorado) recuperable con John the Ripper. & \cellcolor{csfRS}\textbf{9.8 (Crítico)} \\
\hline

\textbf{CWE-256} / \textbf{CWE-312} & \textbf{ACT-09} y \textbf{ACT-08} & Almacenamiento de credenciales en texto claro dentro de respaldos sin cifrar. & Dentro del directorio /var/backups/ se guarda un archivo comprimido de respaldo que incluye un archivo de texto con las contraseñas en claro (developer2024 : LetMeIn2024!) para acceder a codeserver. & \cellcolor{csfDE}\textbf{8.8 (Alto)} \\
\hline

\textbf{CWE-284} (Túnel SSH / Pivoteo) & \textbf{ACT-06} y \textbf{ACT-08} & Falta de segmentación interna y reenvío de puertos habilitado (AllowTcpForwarding yes). & La doble interfaz de red de webserver (192.168.177.129 y 10.0.0.78) permite levantar túneles SSH (ssh -L) desde la DMZ para alcanzar los puertos 3389/tcp (RDP) y 5985/tcp (WinRM) de codeserver. & \cellcolor{csfDE}\textbf{8.6 (Alto)} \\
\hline

\textbf{CVE-2017-9798} (\textit{Optionsbleed}) / CWE-1104 & \textbf{ACT-07} (Apache 2.4.10) & Servidor web obsoleto (Apache/2.4.10) fuera de soporte y con fuga de memoria en cabeceras HTTP. & El servicio HTTP en el puerto 80/tcp muestra su versión en las cabeceras de respuesta y no cuenta con parches de seguridad recientes. & \cellcolor{csfDE}\textbf{7.5 (Alto)} \\
\hline

\textbf{OWASP LLM01 / LLM06} y \textbf{CWE-200} & \textbf{ACT-03}, \textbf{ACT-04} y \textbf{ACT-01} & Inyección de prompts (\textit{Prompt Injection}), fuga de código por \textit{ShadowAI} y envenenamiento de datos. & El personal de desarrollo llega a pegar fragmentos del archivo de Python que es el bot o del dataset en modelos públicos, además de que el bot carece de filtros estrictos en las entradas de los usuarios. & \cellcolor{csfDE}\textbf{8.5 (Alto)} \\
\hline

\textbf{CWE-326} / \textbf{CWE-353} & \textbf{ACT-10} y \textbf{ACT-07} & Aceptación de protocolos antiguos (TLS 1.0/1.1), cifrados CBC y recursos externos sin SRI. & La compatibilidad con versiones viejas de TLS debilita el cifrado del canal y la carga de scripts externos sin validación de integridad (SRI) expone el portal a modificaciones. & \cellcolor{csfGV}\textbf{6.5 (Medio)} \\
\hline
\end{longtable}
\endgroup

% ---------------------------------------------------------
\subsection{Registro de Actores de Amenaza}
% ---------------------------------------------------------

En la Tabla~\ref{tab:actores-amenaza} se describen los actores de amenaza internos y externos identificados para la operación de la empresa, considerando incidentes frecuentes en el sector como el robo de bases de datos mediante vulnerabilidades en servidores web expuestos, cadenas de escalación de privilegios, movimientos laterales hacia sistemas Windows antiguos y fugas de información en plataformas externas de inteligencia artificial.

\vspace{0.2cm}
\begingroup
\scriptsize
\setlength{\tabcolsep}{4pt}
\renewcommand{\arraystretch}{1.28}
\noindent
\begin{longtable}{|C{1.3cm}|L{2.9cm}|L{2.2cm}|L{2.8cm}|L{5.1cm}|C{2.0cm}|}
\caption{Registro de actores de amenaza para \textit{Aetheria FinTech \& AI Solutions}.}
\label{tab:actores-amenaza} \\
\hline
\rowcolor{xlSummary}
\textbf{ID} & \textbf{Actor de Amenaza} & \textbf{Categoría / Origen} & \textbf{Motivación y Recursos} & \textbf{Tácticas, Técnicas y Procedimientos (TTPs)} & \textbf{Activos Objetivo} \\
\hline
\endfirsthead
\hline
\rowcolor{xlSummary}
\textbf{ID} & \textbf{Actor de Amenaza} & \textbf{Categoría / Origen} & \textbf{Motivación y Recursos} & \textbf{Tácticas, Técnicas y Procedimientos (TTPs)} & \textbf{Activos Objetivo} \\
\hline
\endhead
\hline
\endlastfoot

\textbf{THR-01} & \textbf{Grupos de Ransomware y Cibercrimen Financiero} (\textit{LockBit, Cl0p, BlackCat}) & Externo / Crimen Organizado & \textbf{Motivación:} Obtener ganancias mediante extorsión y venta de datos personales.\newline \textbf{Recursos:} Altos. & Escaneo de puertos SSH y HTTP (T1595), fuerza bruta con diccionarios (T1110), escalación local con \textit{PwnKit} (T1068), lectura de contraseñas en respaldos (T1552.001), pivoteo por túnel SSH (T1572) hacia codeserver y cifrado de la base de datos. & \textbf{ACT-01, ACT-02, ACT-06, ACT-08, ACT-09} \\
\hline

\textbf{THR-02} & \textbf{Competidores Desleales y Brokers de Propiedad Intelectual} & Externo / Espionaje Corporativo & \textbf{Motivación:} Robar los pesos del modelo de IA, el archivo de Python que es el bot y el dataset.\newline \textbf{Recursos:} Medios a altos. & Uso de credenciales filtradas (T1078), conexión remota a codeserver por WinRM o RDP (T1021.001, T1021.006), copia del archivo de Python que es el bot y extracción del dataset de entrenamiento. & \textbf{ACT-03, ACT-04, ACT-08, ACT-09} \\
\hline

\textbf{THR-03} & \textbf{Amenaza Interna No Maliciosa (\textit{Insider Negligente / ShadowAI})} & Interno / Desarrolladores y Analistas & \textbf{Motivación:} Resolver errores de programación rápido durante la jornada laboral.\newline \textbf{Recursos:} Cuentas legítimas (developer2024, webadmin). & Guardar un archivo de texto con las contraseñas sin cifrar, dejar respaldos comprimidos en /var/backups/ y pegar partes del archivo de Python que es el bot o del dataset en servicios externos de IA (GroqCloud, Claude, DeepSeek). & \textbf{ACT-01, ACT-03, ACT-04, ACT-09} \\
\hline

\textbf{THR-04} & \textbf{Amenaza Interna Maliciosa (\textit{Empleado Disconforme})} & Interno / Administradores o Exempleados & \textbf{Motivación:} Represalia laboral o beneficio económico por fraude.\newline \textbf{Recursos:} Acceso directo a la red interna 10.0.0.0/24. & Modificación intencional de las reglas de crédito en el archivo de Python que es el bot, alteración de etiquetas en el dataset de entrenamiento (\textit{Data Poisoning}) o borrado de registros en dbserver. & \textbf{ACT-01, ACT-02, ACT-03, ACT-04} \\
\hline

\textbf{THR-05} & \textbf{Atacantes de Sistemas de IA (\textit{Prompt Injectors})} & Externo / Usuarios Maliciosos y Defraudadores & \textbf{Motivación:} Alterar las respuestas del bot para aprobar créditos o extraer datos de la base de datos.\newline \textbf{Recursos:} Bajos a medios. & Envío de instrucciones manipuladas mediante \textit{Prompt Injection} (AML.T0051), salto de restricciones del modelo (\textit{Jailbreak}) y consultas indebidas desde el portal web Apache. & \textbf{ACT-01, ACT-03, ACT-07} \\
\hline

\textbf{THR-06} & \textbf{Grupos Hacktivistas y Botnets de DDoS / Cryptojacking} & Externo / Colectivos y Atacantes Oportunistas & \textbf{Motivación:} Tirar el portal financiero o usar las GPUs para minar criptomonedas.\newline \textbf{Recursos:} Medios. & Ataques de denegación de servicio (DDoS) contra Apache/2.4.10 (T1498), alteración visual del portal web y ejecución de mineros en el clúster de servidores GPU (T1496). & \textbf{ACT-05, ACT-06, ACT-07, ACT-10} \\
\hline
\end{longtable}
\endgroup

% =========================================================
\section{Registro de Riesgos (10 Escenarios de Riesgo)}
% =========================================================

A partir de los activos y amenazas anteriores, se definieron diez escenarios de riesgo (RK-001 a RK-010) para \textbf{Aetheria FinTech \& AI Solutions S.A. de C.V.}, redactados tanto por su impacto operativo (Forma 1) como por el actor y activo involucrado (Forma 2), además de su registro general.

% ---------------------------------------------------------
\subsection{Escenarios de Riesgo en Forma 1 y Forma 2}
% ---------------------------------------------------------

\begingroup
\scriptsize
\setlength{\tabcolsep}{3.5pt}
\renewcommand{\arraystretch}{1.28}
\noindent
\begin{longtable}{|C{0.6cm}|C{1.3cm}|L{2.8cm}|L{5.8cm}|L{5.8cm}|}
\caption{Escenarios de riesgo redactados en Forma 1 y Forma 2.}
\label{tab:escenarios-formas} \\
\hline
\rowcolor{xlSummary}
\textbf{\#} & \textbf{ID} & \textbf{Nombre del Riesgo} & \textbf{Escenario en Forma 1 (Orientado al Impacto y Proceso Operativo)} & \textbf{Escenario en Forma 2 (Orientado al Actor y Activo Comprometido)} \\
\hline
\endfirsthead
\hline
\rowcolor{xlSummary}
\textbf{\#} & \textbf{ID} & \textbf{Nombre del Riesgo} & \textbf{Escenario en Forma 1 (Orientado al Impacto y Proceso Operativo)} & \textbf{Escenario en Forma 2 (Orientado al Actor y Activo Comprometido)} \\
\hline
\endhead
\hline
\endlastfoot

1 & RK-001 & Compromiso de webserver, escalación a root y pivoteo hacia codeserver & Compromiso total del servidor Linux frontal (webserver) y acceso lateral hacia el servidor interno de producción (codeserver) debido a contraseñas débiles en SSH (13337/tcp), falta de parches en pkexec (CVE-2021-4034) y contraseñas guardadas en texto claro dentro de respaldos, provocando el robo o alteración del archivo de Python que es el bot y la caída del servicio. & Un grupo de cibercrimen (\textbf{THR-01}) entra a webserver (\textbf{ACT-06}) por fuerza bruta en SSH, escala privilegios a root con \textit{PwnKit}, obtiene las credenciales de developer2024 del archivo comprimido de respaldo en /var/backups/ (\textbf{ACT-09}) y abre un túnel SSH hacia codeserver (\textbf{ACT-08}) para tomar control del archivo de Python que es el bot (\textbf{ACT-03}). \\
\hline

2 & RK-002 & ShadowAI (Fuga del archivo de Python que es el bot y dataset en IA externa) & Un desarrollador utiliza plataformas externas de inteligencia artificial (GroqCloud, Claude, DeepSeek) sin permiso de la empresa para revisar errores del archivo de Python que es el bot y probar muestras del dataset, ocasionando la filtración de la lógica del modelo, credenciales de base de datos y datos personales de clientes. & Un colaborador interno del área de desarrollo (\textbf{THR-03}) compromete la confidencialidad del archivo de Python que es el bot (\textbf{ACT-03}), del dataset de entrenamiento (\textbf{ACT-04}) y de la base de datos (\textbf{ACT-01}) al subirlos a servidores externos de IA (CWE-200, OWASP LLM06). \\
\hline

3 & RK-003 & Exfiltración y cifrado por Ransomware de la Base de Datos del Modelo de IA & Robo masivo de información y cifrado por ransomware de la Base de Datos del Modelo de IA y registros de clientes (ACT-01) por falta de doble factor de autenticación en la red 10.0.0.0/24 y reutilización de contraseñas en dbserver (ACT-02), generando extorsión económica y multas regulatorias. & Un grupo de ransomware (\textbf{THR-01}), una vez dentro de la subred 10.0.0.0/24, ingresa al servidor de base de datos (\textbf{ACT-02}) con credenciales locales, extrae los 450,000 registros de clientes (\textbf{ACT-01}) y cifra las tablas de producción. \\
\hline

4 & RK-004 & Envenenamiento (\textit{Data Poisoning}) del Dataset de entrenamiento del bot & Alteración maliciosa y sesgo en las evaluaciones crediticias del modelo de IA (ACT-03) a causa de la modificación del dataset de entrenamiento (ACT-04) por no contar con verificación de integridad mediante hashes SHA-256, lo que deriva en la aprobación de préstamos incobrables. & Un actor interno malicioso (\textbf{THR-04}) o intruso altera las etiquetas de riesgo en el dataset de entrenamiento (\textbf{ACT-04}) antes de que se ejecute el reentrenamiento en el clúster GPU (\textbf{ACT-05}), dañando la precisión del modelo (\textbf{ACT-03}). \\
\hline

5 & RK-005 & Manipulación del archivo de Python que es el bot vía \textit{Prompt Injection} & Autorización indebida de créditos y exposición de información de otros clientes mediante el archivo de Python que es el bot (ACT-03) debido a ataques de \textit{Prompt Injection} (OWASP LLM01) y a la falta de validación entre las salidas del modelo y las consultas hacia la base de datos (ACT-01). & Usuarios maliciosos externos (\textbf{THR-05}) introducen instrucciones manipuladas desde el portal web (\textbf{ACT-07}) hacia el archivo de Python que es el bot (\textbf{ACT-03}), saltando las reglas del asistente para consultar saldos ajenos en \textbf{ACT-01} o forzar créditos aprobados. \\
\hline

6 & RK-006 & Caída del portal web por explotación de Apache/2.4.10 o ataque DDoS & Interrupción del portal web y de la conexión con el bot debido al uso de una versión antigua de Apache/2.4.10 (ACT-07) o a un ataque de denegación de servicio distribuido (DDoS) que agota la memoria y el ancho de banda de webserver (ACT-06). & Un grupo hacktivista o competidor (\textbf{THR-06}) lanza un ataque DDoS y explota fallos en Apache/2.4.10 (\textbf{ACT-07}, CVE-2017-9798), dejando sin servicio el portal de la empresa e incumpliendo la disponibilidad del $99.9\%$. \\
\hline

7 & RK-007 & Secuestro de cómputo (\textit{Cryptojacking}) o fallo en la infraestructura física GPU & Pérdida de capacidad de procesamiento e interrupción del entrenamiento del modelo en la infraestructura física de servidores GPU (ACT-05) por la instalación de mineros de criptomonedas (\textit{Cryptojacking}) o por sobrecalentamiento de los equipos sin respaldo activo. & Atacantes oportunistas (\textbf{THR-06}) acceden a los nodos del clúster GPU (\textbf{ACT-05}) mediante paneles de administración mal protegidos e instalan programas de minería que consumen el $100\%$ de las tarjetas gráficas, frenando el entrenamiento del bot. \\
\hline

8 & RK-008 & Compromiso de codeserver por servicios RDP y WinRM en Windows 2012 R2 & Toma de control del servidor interno codeserver (ACT-08) por mantener en operación un sistema operativo fuera de soporte (\textbf{Windows Server 2012 R2 Standard}) con los puertos de administración remota RDP (3389/tcp) y WinRM (5985/tcp) abiertos sin autenticación multifactor. & Un atacante externo o competidor (\textbf{THR-02}) emplea las credenciales de la cuenta developer2024 en los servicios WinRM y RDP de codeserver (\textbf{ACT-08}), ejecutando sesiones remotas con Evil-WinRM para copiar el archivo de Python que es el bot (\textbf{ACT-03}). \\
\hline

9 & RK-009 & Degradación criptográfica TLS (1.0/1.1) e inyección por falta de SRI & Modificación no autorizada del portal web en el navegador del usuario o intercepción de tráfico por permitir versiones obsoletas de TLS (1.0 y 1.1) y cargar librerías externas de JavaScript sin verificación de integridad (SRI) en ACT-07 y ACT-10. & Un atacante externo intercepta la sesión de un cliente degradando el protocolo TLS en \textbf{ACT-10} o altera un recurso externo cargado sin validación de integridad en el servidor Apache (\textbf{ACT-07}) para robar datos de acceso. \\
\hline

10 & RK-010 & Pérdida de archivos de respaldo por almacenamiento local en DMZ & Imposibilidad de recuperar la versión funcional del modelo de IA (ACT-03) y de la plataforma debido a que los archivos comprimidos de respaldo (ACT-09 en /root/ y /var/backups/) están guardados dentro del mismo servidor frontal expuesto (webserver). & Durante una intrusión o daño de almacenamiento en webserver (\textbf{ACT-06}), el atacante (\textbf{THR-01}) borra tanto el sistema activo como los archivos comprimidos de respaldo locales (\textbf{ACT-09}), dejando a la empresa sin copias de recuperación. \\
\hline
\end{longtable}
\endgroup

% ---------------------------------------------------------
\subsection{Plantilla Completa del Registro de Riesgos}
% ---------------------------------------------------------

\begingroup
\tiny
\setlength{\tabcolsep}{2.2pt}
\renewcommand{\arraystretch}{1.35}
\noindent
\begin{longtable}{|C{0.80cm}|L{1.75cm}|L{1.20cm}|L{1.85cm}|L{1.65cm}|L{1.80cm}|L{1.70cm}|C{1.15cm}|C{1.30cm}|C{1.35cm}|C{1.35cm}|}
\caption{Registro general de riesgos de \textit{Aetheria FinTech \& AI Solutions S.A. de C.V.}}
\label{tab:plantilla-registro-riesgos} \\
\hline
\rowcolor{xlSummary}
\textbf{ID} & \textbf{Nombre del Riesgo} & \textbf{Activo Afectado} & \textbf{Vulnerabilidad} & \textbf{Actor de Amenaza} & \textbf{Impacto} & \textbf{Propietario del Riesgo} & \textbf{Nivel de Riesgo} & \textbf{Respuesta al Riesgo} & \textbf{Estado Plan de Acción} & \textbf{Última Evaluación} \\
\hline
\endfirsthead
\hline
\rowcolor{xlSummary}
\textbf{ID} & \textbf{Nombre del Riesgo} & \textbf{Activo Afectado} & \textbf{Vulnerabilidad} & \textbf{Actor de Amenaza} & \textbf{Impacto} & \textbf{Propietario del Riesgo} & \textbf{Nivel de Riesgo} & \textbf{Respuesta al Riesgo} & \textbf{Estado Plan de Acción} & \textbf{Última Evaluación} \\
\hline
\endhead
\hline
\endlastfoot

\textbf{RK-001} & Compromiso de webserver, escalación a root y pivoteo a codeserver & \textbf{ACT-06, ACT-08, ACT-03, ACT-09} & CVE-2021-4034 (\textit{PwnKit}), CWE-521, CWE-256, CWE-284 & \textbf{THR-01} (Ransomware / Cibercrimen) & Pérdida total de \textbf{C, I, D}: control de webserver y codeserver, robo del archivo de Python que es el bot y paro operativo. & CISO y Líder de Infraestructura TI & \cellcolor{csfRS}\textbf{CRÍTICO (16.25)} & \textbf{Mitigar} & En implementación (60\%) & 14/11/2026 \\
\hline

\textbf{RK-002} & ShadowAI (Fuga del archivo de Python que es el bot y dataset) & \textbf{ACT-03, ACT-04, ACT-01} & OWASP LLM06, CWE-200 (Ausencia de DLP y gobierno IA) & \textbf{THR-03} (Insider Negligente / Desarrollo) & Pérdida de \textbf{Confidencialidad}: fuga de propiedad intelectual del modelo y sanciones por exposición de PII. & CISO, DPO y Director de IA & \cellcolor{csfDE}\textbf{ALTO (11.00)} & \textbf{Mitigar / Evitar} & En implementación (70\%) & 14/11/2026 \\
\hline

\textbf{RK-003} & Exfiltración y cifrado por Ransomware de la BD del Modelo de IA & \textbf{ACT-01, ACT-02} & CWE-306, CWE-522 (Credenciales locales sin MFA ni cifrado) & \textbf{THR-01} (Grupos RaaS / Extorsión) & Pérdida severa de \textbf{C, I, D}: secuestro de 450,000 registros financieros, extorsión y multas CNBV. & Director de Base de Datos (DBA) y CISO & \cellcolor{csfRS}\textbf{CRÍTICO (15.00)} & \textbf{Mitigar / Transferir} & En progreso (40\%) & 14/11/2026 \\
\hline

\textbf{RK-004} & Envenena\-miento (\textit{Data Poisoning}) del Dataset de entrenamiento & \textbf{ACT-04, ACT-03, ACT-05} & AML.T0020, CWE-345 (Falta de firma SHA-256 del corpus) & \textbf{THR-04} (Insider Malicioso) / \textbf{THR-02} & Pérdida de \textbf{Integridad}: degradación del modelo predictivo y otorgamiento de créditos incobrables. & Líder de Ciencia de Datos e IA & \cellcolor{csfDE}\textbf{ALTO (10.50)} & \textbf{Mitigar} & Planificado & 14/11/2026 \\
\hline

\textbf{RK-005} & Manipulación del archivo de Python que es el bot vía \textit{Prompt Injection} & \textbf{ACT-03, ACT-01, ACT-07} & OWASP LLM01, CWE-77 (Falta de guardrails y privilegios BD) & \textbf{THR-05} (Defraudadores y Atacantes de IA) & Pérdida de \textbf{Integridad y Confidencialidad}: préstamos fraudulentos autorizados por el bot y fuga de saldos. & Líder de Desarrollo de IA & \cellcolor{csfDE}\textbf{ALTO (10.00)} & \textbf{Mitigar} & En pruebas (50\%) & 14/11/2026 \\
\hline

\textbf{RK-006} & Caída del portal financiero por explotación de Apache/2.4.10 o DDoS & \textbf{ACT-07, ACT-06, ACT-10} & CVE-2017-9798 (\textit{Optionsbleed}), CWE-400 (Servidor web EoL) & \textbf{THR-06} (Hacktivismo / Botnets DDoS) & Pérdida de \textbf{Disponibilidad}: interrupción del portal web y API del bot, penalización de SLA. & Líder de Redes y Operaciones & \cellcolor{csfDE}\textbf{ALTO (9.75)} & \textbf{Mitigar} & En progreso (65\%) & 14/11/2026 \\
\hline

\textbf{RK-007} & Secuestro de cómputo (\textit{Cryptojacking}) o fallo en clúster GPU & \textbf{ACT-05} & CWE-250, contenedores sin cuotas y falta de failover GPU & \textbf{THR-06} (Oportunistas) / Amenaza Ambiental & Pérdida de \textbf{Disponibilidad}: retraso en reentrenamiento del bot, sobrecosto eléctrico y daño de GPUs. & Gerente de Centro de Datos & \cellcolor{csfGV}\textbf{MEDIO (6.75)} & \textbf{Mitigar / Transferir} & Planificado & 14/11/2026 \\
\hline

\textbf{RK-008} & Compromiso de codeserver por RDP/WinRM en Windows 2012 R2 & \textbf{ACT-08, ACT-03} & CWE-1104 (SO sin soporte EoL), CWE-308 (WinRM 5985 y RDP 3389 sin MFA) & \textbf{THR-02} (Espionaje Corporativo) / \textbf{THR-01} & Pérdida de \textbf{Confidencialidad e Integridad}: ejecución remota de comandos y robo del archivo de Python que es el bot. & Líder de Sistemas Windows & \cellcolor{csfDE}\textbf{ALTO (12.00)} & \textbf{Mitigar} & En ejecución (45\%) & 14/11/2026 \\
\hline

\textbf{RK-009} & Interceptación o inyección web por TLS legado (1.0/1.1) y falta de SRI & \textbf{ACT-10, ACT-07} & CWE-326 (TLS débil / CBC), CWE-353 (Ausencia de validación SRI) & \textbf{THR-01} / Atacantes \textit{Man-in-the-Middle} & Pérdida de \textbf{Confidencialidad e Integridad}: captura de credenciales de clientes y alteración del frontend web. & Líder de Seguridad de Aplicaciones & \cellcolor{csfGV}\textbf{MEDIO (5.50)} & \textbf{Mitigar} & Completado (80\%) & 14/11/2026 \\
\hline

\textbf{RK-010} & Pérdida de archivos de respaldo por almacenamiento local en DMZ & \textbf{ACT-09, ACT-06} & CWE-312, ausencia de regla de respaldo 3-2-1 e inmutabilidad & \textbf{THR-01} (Ransomware) / Fallo de disco & Pérdida de \textbf{Disponibilidad y Confidencialidad}: imposibilidad de recuperar el sistema tras un ataque. & Director de Continuidad (DRP) & \cellcolor{csfDE}\textbf{ALTO (9.00)} & \textbf{Mitigar} & En implementación (50\%) & 14/11/2026 \\
\hline
\end{longtable}
\endgroup

\newpage
% =========================================================
\section{Análisis Detallado de Dos Riesgos Relevantes y Planes de Respuesta}
% =========================================================

A continuación se presenta la evaluación cuantitativa, el plan de respuesta y los indicadores de monitoreo para dos de los riesgos principales de la organización:
\begin{enumerate}[noitemsep,topsep=2pt]
    \item \textbf{Caso 1 (RK-001):} Compromiso del servidor Linux frontal (webserver), escalación de privilegios a root (CVE-2021-4034), obtención del archivo de texto con las contraseñas en los respaldos y pivoteo por túnel SSH hacia codeserver (10.0.0.201) para tomar control del archivo de Python que es el bot.
    \item \textbf{Caso 2 (RK-002):} Fuga del archivo de Python que es el bot, de muestras del dataset de entrenamiento y de esquemas de la base de datos por el uso no autorizado de servicios externos de IA (ShadowAI).
\end{enumerate}

% ---------------------------------------------------------
\subsection{Evaluación del Riesgo 1 (RK-001): Intrusión en webserver, Escalación a root y Pivoteo hacia codeserver}
% ---------------------------------------------------------

\begingroup
\scriptsize
\setlength{\tabcolsep}{3.5pt}
\renewcommand{\arraystretch}{1.26}
\noindent
\begin{longtable}{|L{3.5cm}|C{2.0cm}|C{2.0cm}|C{2.0cm}|C{2.0cm}|C{2.0cm}|C{2.0cm}|}
\hline
% ==================== SECCIÓN I - RESUMEN ====================
\multicolumn{7}{|c|}{\cellcolor{xlHeader}\textbf{Sección I - Resumen (RK-001)}} \\
\hline
\textbf{ID del riesgo} & \multicolumn{3}{l|}{RK-001} & \multicolumn{2}{l|}{\textbf{Nivel de Riesgo}} & \cellcolor{xlSummary}\textbf{CRÍTICO (16.25)} \\
\hline
\textbf{Nombre del riesgo} & \multicolumn{3}{p{6.0cm}|}{Compromiso de webserver, escalación a root (CVE-2021-4034) y pivoteo SSH hacia codeserver (el archivo de Python que es el bot)} & \multicolumn{2}{l|}{\textbf{Respuesta al riesgo}} & \cellcolor{xlSummary}\textbf{MITIGAR} \\
\hline
\textbf{Propietario del riesgo} & \multicolumn{3}{p{6.0cm}|}{CISO y Líder de Infraestructura de TI} & \multicolumn{2}{l|}{\textbf{Estado del plan de acción}} & \cellcolor{xlSummary}\textbf{EN IMPLEMENTACIÓN (60\%)} \\
\hline
\textbf{Descripción del riesgo} & \multicolumn{3}{p{6.0cm}|}{Intrusión por fuerza bruta SSH en webserver, escalación con \textit{PwnKit}, robo del archivo de texto con las contraseñas y túnel SSH hacia codeserver} & \multicolumn{2}{l|}{\textbf{Última evaluación del riesgo}} & \cellcolor{xlSummary}\textbf{14/11/2026} \\
\hline

% ==================== SECCIÓN II - DESCRIPCIÓN DEL RIESGO ====================
\multicolumn{7}{|c|}{\cellcolor{xlHeader}\textbf{Sección II - Descripción del riesgo (RK-001)}} \\
\hline
\textbf{Nombre del riesgo} & \multicolumn{6}{p{12.8cm}|}{Compromiso del servidor Linux frontal (webserver), escalación de privilegios a root y movimiento lateral hacia el servidor interno Windows (codeserver) donde se encuentra el archivo de Python que es el bot} \\
\hline
\textbf{Categoría de riesgo} & Estratégico & Programa & Proyecto & \cellcolor{xlSelected}\textbf{Operativo} & Otro: & \\
\hline
\textbf{Escenario de alto nivel} & \multicolumn{6}{p{12.8cm}|}{Un atacante externo localiza mediante un escaneo de puertos el servicio OpenSSH 6.7p1 en el puerto 13337/tcp y Apache/2.4.10 en el puerto 80/tcp del servidor webserver (192.168.177.129). Obtiene acceso a la cuenta webadmin por fuerza bruta con Hydra (contraseña kitty123), escala privilegios a root explotando CVE-2021-4034 (\textit{PwnKit}) en /usr/bin/pkexec, extrae del archivo comprimido de respaldo en /var/backups/ el archivo de texto con las contraseñas (developer2024 : LetMeIn2024!) y usa la segunda interfaz de red (10.0.0.78) para crear un túnel SSH hacia codeserver (10.0.0.201), logrando acceder por RDP o WinRM al archivo de Python que es el bot.} \\
\hline
\textbf{Descripción detallada del riesgo} & \textbf{Evento} & \multicolumn{5}{p{10.5cm}|}{Toma de control de webserver en la DMZ y salto hacia codeserver en la red interna 10.0.0.0/24, comprometiendo el archivo de Python que es el bot} \\
\cline{2-7}
 & \textbf{Activo / Recurso} & \multicolumn{5}{p{10.5cm}|}{webserver (\textbf{ACT-06}), Apache/2.4.10 (\textbf{ACT-07}), archivos comprimidos de respaldo (\textbf{ACT-09}), codeserver (\textbf{ACT-08}) y el archivo de Python que es el bot (\textbf{ACT-03})} \\
\cline{2-7}
 & \textbf{Tipo de Amenaza} & \multicolumn{5}{p{10.5cm}|}{\cellcolor{xlSummary}Intrusión Externa, Escalación Local de Privilegios (CVE-2021-4034), Exposición de Credenciales (CWE-256) y Pivoteo de Red (CWE-284)} \\
\cline{2-7}
 & \textbf{Actor} & \multicolumn{5}{p{10.5cm}|}{\cellcolor{xlSummary}\textbf{THR-01} (Grupos de Ransomware y Cibercrimen Financiero) y \textbf{THR-02} (Competidores / Brokers de Acceso Inicial)} \\
\cline{2-7}
 & \textbf{Probabilidad} & \multicolumn{5}{p{10.5cm}|}{Muy Alta} \\
\cline{2-7}
 & \textbf{Consecuencias} & \multicolumn{5}{p{10.5cm}|}{\cellcolor{xlComment}Robo o modificación del archivo de Python que es el bot, pérdida de propiedad intelectual, caída del portal web e incumplimiento del nivel de servicio con clientes financieros} \\
\hline
\multicolumn{2}{|l|}{\textbf{Otra información del escenario}} & \multicolumn{5}{p{10.5cm}|}{El ataque combina el escaneo de puertos, ataque de diccionario sobre SSH, uso de exploit local en Polkit, lectura de respaldos sin cifrar y reenvío de puertos SSH hacia el segmento interno.} \\
\hline

% ==================== SECCIÓN III - ANÁLISIS DE RIESGO ====================
\multicolumn{7}{|c|}{\cellcolor{xlHeader}\textbf{Sección III - Análisis de Riesgo (RK-001)}} \\
\hline
\multicolumn{7}{|c|}{\textbf{Frecuencia}} \\
\hline
\textbf{Ocurrencia Anual} & 0.0 & 1.0 & 2.0 & 3.0 & 4.0 & \cellcolor{xlSelected}\textbf{5.0} \\
\cline{2-7}
 & $N \leq 0.01$ & $0.01 < N \leq 0.1$ & $0.1 < N \leq 1$ & $1 < N \leq 10$ & $10 < N \leq 100$ & \cellcolor{xlSelected}\boldmath$100 < N$ \\
\hline
\textbf{Comentarios sobre la frecuencia} & \multicolumn{6}{p{12.8cm}|}{\cellcolor{xlComment}Al estar expuesto directamente en la DMZ, el servidor webserver recibe de forma constante más de 100 intentos de escaneo y ataques de fuerza bruta por SSH al año ($100 < N$), además de que el exploit para CVE-2021-4034 es público y muy fácil de ejecutar.} \\
\hline
\multicolumn{7}{|c|}{\textbf{Impacto}} \\
\hline
\textbf{Pérdida de ganancia anual} & 0.0 & 1.0 & 2.0 & 3.0 & \cellcolor{xlSelected}\textbf{4.0} & 5.0 \\
\cline{2-7}
 & $I \leq 0.1\%$ & $0.1\% < I \leq 1\%$ & $1\% < I \leq 3\%$ & $3\% < I \leq 5\%$ & \cellcolor{xlSelected}\boldmath$5\% < I \leq 10\%$ & $10\% < I$ \\
\hline
\textbf{Comentarios sobre la pérdida} & \multicolumn{6}{p{12.8cm}|}{\cellcolor{xlComment}Un paro en el portal Apache y la alteración del archivo de Python que es el bot en codeserver detendrían el cobro por consultas API y la evaluación de créditos, causando pérdidas estimadas entre el 5\% y el 10\% de los ingresos anuales.} \\
\hline
\textbf{Costo de la respuesta al riesgo} & 0.0 & 1.0 & \cellcolor{xlSelected}\textbf{2.0} & 3.0 & 4.0 & 5.0 \\
\cline{2-7}
 & $I \leq \$10\text{k}$ & $\$10\text{k} < I \leq \$100\text{k}$ & \cellcolor{xlSelected}\boldmath$\$100\text{k} < I \leq \$1\text{M}$ & $\$1\text{M} < I \leq \$10\text{M}$ & $\$10\text{M} < I \leq \$100\text{M}$ & $\$100\text{M} < I$ \\
\hline
\textbf{Comentarios sobre el costo de la respuesta} & \multicolumn{6}{p{12.8cm}|}{Actualizar webserver (de Debian 8 a Debian 12) y codeserver (de Windows Server 2012 R2 a Windows Server 2022), junto con la puesta en marcha de un firewall interno y un gestor de secretos, tiene un costo aproximado de \$180,000 USD ($\$100\text{k} < I \leq \$1\text{M}$).} \\
\hline
\textbf{Impacto reputacional} & 0.0 & 1.0 & 2.0 & 3.0 & \cellcolor{xlSelected}\textbf{4.0} & 5.0 \\
\cline{2-7}
 & $I \leq 0.5$ & $0.5 < I \leq 1$ & $1 < I \leq 1.5$ & $1.5 < I \leq 2$ & \cellcolor{xlSelected}\boldmath$2 < I \leq 2.5$ & $2.5 < I$ \\
\hline
\textbf{Comentarios sobre el impacto reputacional} & \multicolumn{6}{p{12.8cm}|}{La afectación del archivo de Python que es el bot en producción y la interrupción del servicio generarían desconfianza directa entre las financieras y cajas de ahorro clientes, con un impacto de entre 2.0 y 2.5 puntos.} \\
\hline
\textbf{Costo de multas por incumplimiento} & 0.0 & 1.0 & 2.0 & \cellcolor{xlSelected}\textbf{3.0} & 4.0 & 5.0 \\
\cline{2-7}
 & $\$0.00$ & $< \$1\text{M}$ & $< \$10\text{M}$ & \cellcolor{xlSelected}\boldmath$< \$100\text{M}$ & $< \$1\text{B}$ & $> \$1\text{B}$ \\
\hline
\textbf{Comentarios sobre el incumplimiento regulatorio} & \multicolumn{6}{p{12.8cm}|}{\cellcolor{xlComment}Operar servidores con sistemas fuera de soporte (Debian 8 y Windows Server 2012 R2) y guardar contraseñas en texto claro incumple los controles exigidos por la CNBV y la LFPDPPP, lo que puede derivar en multas y penalizaciones contractuales ($< \$100\text{M}$).} \\
\hline
\multicolumn{4}{|l|}{\textbf{Calificación general del impacto (promedio de los 4 aspectos)}} & \multicolumn{3}{r|}{\cellcolor{xlSummary}\textbf{3.25}} \\
\hline
\multicolumn{4}{|l|}{\textbf{Calificación general del nivel de riesgo (Frecuencia $\times$ Impacto)}} & \multicolumn{3}{r|}{\cellcolor{xlSummary}\textbf{16.25}} \\
\hline
\end{longtable}
\endgroup

\vspace{0.2cm}

\begingroup
\footnotesize
\setlength{\tabcolsep}{4pt}
\renewcommand{\arraystretch}{1.32}
\noindent
\begin{longtable}{|C{0.8cm}|L{3.1cm}|C{1.9cm}|C{1.9cm}|C{1.9cm}|C{2.2cm}|C{2.2cm}|}
\hline
% ==================== SECCIÓN IV - RESPUESTA AL RIESGO ====================
\multicolumn{7}{|c|}{\cellcolor{xlHeader}\textbf{Sección IV - Respuesta al riesgo (RK-001)}} \\
\hline
\multicolumn{3}{|l|}{\textbf{Respuesta seleccionada:}} & Aceptar & Transferir & \cellcolor{xlSelected}\textbf{Mitigar} & Evitar \\
\hline
\multicolumn{2}{|c|}{\textbf{Justificación}} & \multicolumn{5}{p{10.9cm}|}{\cellcolor{xlComment}El riesgo tiene un nivel crítico de \textbf{16.25} ($F=5.0 \times I_{\text{prom}}=3.25$) y compromete tanto el servidor perimetral como el servidor interno donde corre el bot. No es posible aceptarlo por el daño operativo ni evitarlo apagando el portal público, por lo que se opta por \textbf{MITIGAR} el riesgo reforzando el acceso SSH, corrigiendo los permisos de pkexec, actualizando los sistemas operativos, borrando el archivo de texto con las contraseñas de los respaldos y bloqueando el reenvío de puertos hacia la red interna 10.0.0.0/24.} \\
\hline
\multicolumn{3}{|l|}{\textbf{Estatus general del plan de acción}} & \multicolumn{4}{p{8.8cm}|}{En implementación (60\% de avance: ya se configuró el acceso por llave en SSH, se retiró el bit SUID en pkexec y se eliminó el archivo de texto con las contraseñas)} \\
\hline
\multicolumn{3}{|l|}{\textbf{Problemas relevantes de las acciones de respuesta}} & \multicolumn{4}{p{8.8cm}|}{Algunas librerías de 32 bits en Debian 8.11 requieren actualización previa y se necesita programar una ventana de mantenimiento nocturna para migrar Windows Server 2012 R2 en codeserver sin afectar las consultas hacia el archivo de Python que es el bot.} \\
\hline
\multicolumn{3}{|l|}{\textbf{Hitos conseguidos en las acciones de respuesta}} & \multicolumn{4}{p{8.8cm}|}{Autenticación por llave pública Ed25519 habilitada en SSH (13337/tcp); reenvío de puertos desactivado (AllowTcpForwarding no); archivo de texto con las contraseñas borrado de los respaldos en /var/backups/ y cambio de claves completado.} \\
\hline

% ==================== PLAN DE ACCIÓN ====================
\multicolumn{7}{|c|}{\cellcolor{xlHeader}\textbf{Plan de acción (RK-001)}} \\
\hline
\textbf{\#} & \multicolumn{4}{c|}{\textbf{Acción de respuesta}} & \textbf{Fecha compromiso} & \textbf{Estatus} \\
\hline
1.0 & \multicolumn{4}{p{9.2cm}|}{Desactivar el inicio de sesión por contraseña (PasswordAuthentication no), el acceso directo de root y el reenvío de puertos (AllowTcpForwarding no) en OpenSSH (13337/tcp) dentro de webserver, usando llaves Ed25519 y Fail2Ban} & 01/11/2026 & Completado \\
\hline
2.0 & \multicolumn{4}{p{9.2cm}|}{Quitar el permiso SUID en /usr/bin/pkexec (CVE-2021-4034) y ntfs-3g (CVE-2017-0358), y actualizar webserver de Debian 8.11 a Debian 12 LTS con Apache 2.4.62} & 08/11/2026 & Completado \\
\hline
3.0 & \multicolumn{4}{p{9.2cm}|}{Borrar el archivo de texto con las contraseñas de los respaldos comprimidos en /var/backups/ y /root/, cambiar la clave de developer2024 y guardar las credenciales en un gestor de secretos con respaldos cifrados fuera del servidor} & 12/11/2026 & Completado \\
\hline
4.0 & \multicolumn{4}{p{9.2cm}|}{Colocar reglas de firewall interno entre la interfaz eth1 (10.0.0.78) de webserver y la subred 10.0.0.0/24 para impedir conexiones desde la DMZ hacia los puertos RDP (3389/tcp) y WinRM (5985/tcp) de codeserver} & 20/11/2026 & En progreso \\
\hline
5.0 & \multicolumn{4}{p{9.2cm}|}{Migrar codeserver (10.0.0.201) de Windows Server 2012 R2 a Windows Server 2022 y permitir la administración del archivo de Python que es el bot únicamente desde un servidor bastión con MFA} & 30/11/2026 & En progreso \\
\hline

% ==================== SECCIÓN V - MONITOREO ====================
\multicolumn{7}{|c|}{\cellcolor{xlHeader}\textbf{Sección V - Monitoreo (RK-001)}} \\
\hline
\multicolumn{2}{|l|}{\textbf{Última evaluación del riesgo}} & \multicolumn{2}{c|}{14/11/2026} & \multicolumn{2}{l|}{\textbf{Siguiente evaluación}} & 14/12/2026 \\
\hline
\multicolumn{7}{|c|}{\textbf{Indicadores clave de riesgo (KRIs)}} \\
\hline
\textbf{\#} & \multicolumn{4}{c|}{\textbf{Indicador}} & \textbf{Medición} & \textbf{Fecha} \\
\hline
1.0 & \multicolumn{4}{p{9.2cm}|}{Intentos de fuerza bruta SSH en el puerto 13337/tcp de webserver bloqueados por Fail2Ban} & 0 accesos exitosos (42 IPs bloqueadas) & 14/11/2026 \\
\hline
2.0 & \multicolumn{4}{p{9.2cm}|}{Cantidad de binarios SUID vulnerables a escalación local (CVE-2021-4034, CVE-2016-5195) o archivos de texto con contraseñas en claro en webserver} & 0 vulnerabilidades / 0 archivos en claro & 14/11/2026 \\
\hline
3.0 & \multicolumn{4}{p{9.2cm}|}{Intentos de conexión TCP desde la interfaz interna de la DMZ (10.0.0.78) hacia los puertos 3389/tcp (RDP) y 5985/tcp (WinRM) de codeserver} & 0 conexiones permitidas & 14/11/2026 \\
\hline
\end{longtable}
\endgroup

\newpage
% ---------------------------------------------------------
\subsection{Evaluación del Riesgo 2 (RK-002): Exfiltración del Archivo de Python que es el Bot, Dataset y Base de Datos por Uso de ShadowAI}
% ---------------------------------------------------------

\begingroup
\scriptsize
\setlength{\tabcolsep}{3.5pt}
\renewcommand{\arraystretch}{1.26}
\noindent
\begin{longtable}{|L{3.5cm}|C{2.0cm}|C{2.0cm}|C{2.0cm}|C{2.0cm}|C{2.0cm}|C{2.0cm}|}
\hline
% ==================== SECCIÓN I - RESUMEN ====================
\multicolumn{7}{|c|}{\cellcolor{xlHeader}\textbf{Sección I - Resumen (RK-002)}} \\
\hline
\textbf{ID del riesgo} & \multicolumn{3}{l|}{RK-002} & \multicolumn{2}{l|}{\textbf{Nivel de Riesgo}} & \cellcolor{xlSummary}\textbf{ALTO (11.00)} \\
\hline
\textbf{Nombre del riesgo} & \multicolumn{3}{p{6.0cm}|}{ShadowAI: Exfiltración del archivo de Python que es el bot, dataset de entrenamiento y esquemas de BD en servicios externos de IA} & \multicolumn{2}{l|}{\textbf{Respuesta al riesgo}} & \cellcolor{xlSummary}\textbf{MITIGAR / EVITAR} \\
\hline
\textbf{Propietario del riesgo} & \multicolumn{3}{p{6.0cm}|}{CISO, Oficial de Privacidad (DPO) y Director de IA} & \multicolumn{2}{l|}{\textbf{Estado del plan de acción}} & \cellcolor{xlSummary}\textbf{EN IMPLEMENTACIÓN (70\%)} \\
\hline
\textbf{Descripción del riesgo} & \multicolumn{3}{p{6.0cm}|}{Filtración del archivo de Python que es el bot, muestras del dataset de entrenamiento y credenciales de BD en plataformas públicas de IA} & \multicolumn{2}{l|}{\textbf{Última evaluación del riesgo}} & \cellcolor{xlSummary}\textbf{14/11/2026} \\
\hline

% ==================== SECCIÓN II - DESCRIPCIÓN DEL RIESGO ====================
\multicolumn{7}{|c|}{\cellcolor{xlHeader}\textbf{Sección II - Descripción del riesgo (RK-002)}} \\
\hline
\textbf{Nombre del riesgo} & \multicolumn{6}{p{12.8cm}|}{Uso no autorizado de plataformas externas de IA (GroqCloud, Claude, DeepSeek) para depurar el archivo de Python que es el bot y revisar muestras del dataset de entrenamiento} \\
\hline
\textbf{Categoría de riesgo} & Estratégico & Programa & Proyecto & \cellcolor{xlSelected}\textbf{Operativo} & Otro: & \\
\hline
\textbf{Escenario de alto nivel} & \multicolumn{6}{p{12.8cm}|}{Un programador o analista con acceso al servidor interno codeserver (10.0.0.201) copia partes del archivo de Python que es el bot, cadenas de conexión a la base de datos del modelo (ACT-01) y registros reales del dataset de entrenamiento (ACT-04) en asistentes externos (como GroqCloud, Claude o DeepSeek) para corregir errores de código más rápido. Al hacer esto sin autorización, el código propietario y los datos personales de los clientes quedan almacenados en servidores externos fuera del control de la empresa.} \\
\hline
\textbf{Descripción detallada del riesgo} & \textbf{Evento} & \multicolumn{5}{p{10.5cm}|}{Envío y exposición del archivo de Python que es el bot, muestras del dataset de entrenamiento y datos de conexión de la BD en servicios externos de IA} \\
\cline{2-7}
 & \textbf{Activo / Recurso} & \multicolumn{5}{p{10.5cm}|}{Modelo de IA y el archivo de Python que es el bot (\textbf{ACT-03}), Dataset para el entrenamiento del bot (\textbf{ACT-04}), Base de datos con el modelo de IA (\textbf{ACT-01}) y Servidor codeserver (\textbf{ACT-08})} \\
\cline{2-7}
 & \textbf{Tipo de Amenaza} & \multicolumn{5}{p{10.5cm}|}{\cellcolor{xlSummary}Amenaza Interna No Maliciosa (\textit{Insider} / Colaborador) y Fuga de Información en Modelos LLM Externos (CWE-200, OWASP LLM06)} \\
\cline{2-7}
 & \textbf{Actor} & \multicolumn{5}{p{10.5cm}|}{\cellcolor{xlSummary}\textbf{THR-03} (Colaborador interno con cuenta de desarrollo developer2024) o atacante externo que compromete cuentas en servicios de terceros} \\
\cline{2-7}
 & \textbf{Probabilidad} & \multicolumn{5}{p{10.5cm}|}{Alta} \\
\cline{2-7}
 & \textbf{Consecuencias} & \multicolumn{5}{p{10.5cm}|}{\cellcolor{xlComment}Pérdida de ventaja competitiva por filtración de las reglas de negocio del archivo de Python que es el bot, exposición del dataset de entrenamiento y sanciones legales por compartir datos personales con terceros} \\
\hline
\multicolumn{2}{|l|}{\textbf{Otra información del escenario}} & \multicolumn{5}{p{10.5cm}|}{Debilidad asociada a CWE-200 y OWASP LLM06 (\textit{Sensitive Information Disclosure}) por la falta de filtros de salida de datos (DLP) y políticas sobre el uso de herramientas de IA.} \\
\hline

% ==================== SECCIÓN III - ANÁLISIS DE RIESGO ====================
\multicolumn{7}{|c|}{\cellcolor{xlHeader}\textbf{Sección III - Análisis de Riesgo (RK-002)}} \\
\hline
\multicolumn{7}{|c|}{\textbf{Frecuencia}} \\
\hline
\textbf{Ocurrencia Anual} & 0.0 & 1.0 & 2.0 & 3.0 & \cellcolor{xlSelected}\textbf{4.0} & 5.0 \\
\cline{2-7}
 & $N \leq 0.01$ & $0.01 < N \leq 0.1$ & $0.1 < N \leq 1$ & $1 < N \leq 10$ & \cellcolor{xlSelected}\boldmath$10 < N \leq 100$ & $100 < N$ \\
\hline
\textbf{Comentarios sobre la frecuencia} & \multicolumn{6}{p{12.8cm}|}{\cellcolor{xlComment}Durante las tareas cotidianas de programación y ajuste del modelo en codeserver, el uso de asistentes web externos para revisar funciones del archivo de Python que es el bot ocasiona entre 10 y 100 eventos de fuga de fragmentos de código o datos al año ($10 < N \leq 100$).} \\
\hline
\multicolumn{7}{|c|}{\textbf{Impacto}} \\
\hline
\textbf{Pérdida de ganancia anual} & 0.0 & 1.0 & 2.0 & 3.0 & \cellcolor{xlSelected}\textbf{4.0} & 5.0 \\
\cline{2-7}
 & $I \leq 0.1\%$ & $0.1\% < I \leq 1\%$ & $1\% < I \leq 3\%$ & $3\% < I \leq 5\%$ & \cellcolor{xlSelected}\boldmath$5\% < I \leq 10\%$ & $10\% < I$ \\
\hline
\textbf{Comentarios sobre la pérdida} & \multicolumn{6}{p{12.8cm}|}{\cellcolor{xlComment}Si la lógica del archivo de Python que es el bot o el dataset financiero quedan expuestos, otros competidores podrían replicar el modelo y se cancelarían contratos empresariales, representando una pérdida de entre el 5\% y el 10\% de la ganancia anual.} \\
\hline
\textbf{Costo de la respuesta al riesgo} & 0.0 & \cellcolor{xlSelected}\textbf{1.0} & 2.0 & 3.0 & 4.0 & 5.0 \\
\cline{2-7}
 & $I \leq \$10\text{k}$ & \cellcolor{xlSelected}\boldmath$\$10\text{k} < I \leq \$100\text{k}$ & $\$100\text{k} < I \leq \$1\text{M}$ & $\$1\text{M} < I \leq \$10\text{M}$ & $\$10\text{M} < I \leq \$100\text{M}$ & $\$100\text{M} < I$ \\
\hline
\textbf{Comentarios sobre el costo de la respuesta} & \multicolumn{6}{p{12.8cm}|}{Se puede aprovechar la infraestructura física de servidores GPU que ya tiene la empresa (\textbf{ACT-05}) para montar un asistente de programación local, por lo que solo se requiere adquirir licencias de DLP y filtrado web con un costo anual cercano a \$40,000 USD ($\$10\text{k} < I \leq \$100\text{k}$).} \\
\hline
\textbf{Impacto reputacional} & 0.0 & 1.0 & 2.0 & \cellcolor{xlSelected}\textbf{3.0} & 4.0 & 5.0 \\
\cline{2-7}
 & $I \leq 0.5$ & $0.5 < I \leq 1$ & $1 < I \leq 1.5$ & \cellcolor{xlSelected}\boldmath$1.5 < I \leq 2$ & $2 < I \leq 2.5$ & $2.5 < I$ \\
\hline
\textbf{Comentarios sobre el impacto reputacional} & \multicolumn{6}{p{12.8cm}|}{En caso de darse a conocer que datos del dataset crediticio o partes del archivo de Python que es el bot se subieron a plataformas públicas, la empresa tendría una afectación reputacional de entre 1.5 y 2.0 puntos frente a sus clientes.} \\
\hline
\textbf{Costo de multas por incumplimiento} & 0.0 & 1.0 & 2.0 & \cellcolor{xlSelected}\textbf{3.0} & 4.0 & 5.0 \\
\cline{2-7}
 & $\$0.00$ & $< \$1\text{M}$ & $< \$10\text{M}$ & \cellcolor{xlSelected}\boldmath$< \$100\text{M}$ & $< \$1\text{B}$ & $> \$1\text{B}$ \\
\hline
\textbf{Comentarios sobre el incumplimiento regulatorio} & \multicolumn{6}{p{12.8cm}|}{\cellcolor{xlComment}Enviar datos personales y crediticios contenidos en el dataset (\textbf{ACT-04}) o en la base de datos (\textbf{ACT-01}) hacia servicios externos sin autorización de los titulares viola la LFPDPPP y las disposiciones financieras aplicables ($< \$100\text{M}$).} \\
\hline
\multicolumn{4}{|l|}{\textbf{Calificación general del impacto (promedio de los 4 aspectos)}} & \multicolumn{3}{r|}{\cellcolor{xlSummary}\textbf{2.75}} \\
\hline
\multicolumn{4}{|l|}{\textbf{Calificación general del nivel de riesgo (Frecuencia $\times$ Impacto)}} & \multicolumn{3}{r|}{\cellcolor{xlSummary}\textbf{11.00}} \\
\hline
\end{longtable}
\endgroup

\vspace{0.2cm}

\begingroup
\footnotesize
\setlength{\tabcolsep}{4pt}
\renewcommand{\arraystretch}{1.32}
\noindent
\begin{longtable}{|C{0.8cm}|L{3.1cm}|C{1.9cm}|C{1.9cm}|C{1.9cm}|C{2.2cm}|C{2.2cm}|}
\hline
% ==================== SECCIÓN IV - RESPUESTA AL RIESGO ====================
\multicolumn{7}{|c|}{\cellcolor{xlHeader}\textbf{Sección IV - Respuesta al riesgo (RK-002)}} \\
\hline
\multicolumn{3}{|l|}{\textbf{Respuesta seleccionada:}} & Aceptar & Transferir & \cellcolor{xlSelected}\textbf{Mitigar} & \cellcolor{xlSelected}\textbf{Evitar} \\
\hline
\multicolumn{2}{|c|}{\textbf{Justificación}} & \multicolumn{5}{p{10.9cm}|}{\cellcolor{xlComment}El riesgo alcanza un nivel alto (\textbf{11.00}) y compromete directamente los activos principales del negocio (el modelo de IA, el archivo de Python que es el bot, el dataset de entrenamiento y la base de datos). Como el costo de implementación es moderado frente a las pérdidas potenciales, se decide \textbf{MITIGAR} el riesgo mediante reglas de DLP, anonimización del dataset y la instalación de un modelo de apoyo local en el clúster GPU (\textbf{ACT-05}), además de \textbf{EVITAR} el uso de plataformas externas bloqueando su acceso por firewall y políticas de grupo.} \\
\hline
\multicolumn{3}{|l|}{\textbf{Estatus general del plan de acción}} & \multicolumn{4}{p{8.8cm}|}{En implementación (70\% de avance: bloqueo de dominios de IA externos y reglas DLP activas en codeserver y equipos de desarrollo)} \\
\hline
\multicolumn{3}{|l|}{\textbf{Problemas relevantes de las acciones de respuesta}} & \multicolumn{4}{p{8.8cm}|}{Se presentaron algunos falsos positivos iniciales en las reglas DLP al revisar código genérico de Python y fue necesario repartir la memoria VRAM del clúster GPU (\textbf{ACT-05}) para correr el asistente local sin retrasar los entrenamientos.} \\
\hline
\multicolumn{3}{|l|}{\textbf{Hitos conseguidos en las acciones de respuesta}} & \multicolumn{4}{p{8.8cm}|}{Bloqueo DNS y reglas DLP funcionando para sitios públicos de IA; política interna de uso de IA firmada por el equipo; enmascaramiento de datos personales aplicado sobre el dataset de entrenamiento (\textbf{ACT-04}).} \\
\hline

% ==================== PLAN DE ACCIÓN ====================
\multicolumn{7}{|c|}{\cellcolor{xlHeader}\textbf{Plan de acción (RK-002)}} \\
\hline
\textbf{\#} & \multicolumn{4}{c|}{\textbf{Acción de respuesta}} & \textbf{Fecha compromiso} & \textbf{Estatus} \\
\hline
1.0 & \multicolumn{4}{p{9.2cm}|}{Configurar políticas de grupo (GPO) en Windows Server (codeserver) y en los equipos de desarrollo para bloquear navegadores con IA integrada y extensiones no autorizadas} & 02/11/2026 & Completado \\
\hline
2.0 & \multicolumn{4}{p{9.2cm}|}{Activar reglas de Endpoint DLP y filtrado web para impedir que se suba o pegue el archivo de Python que es el bot, archivos de datos estructurados del dataset o cadenas de conexión en sitios externos (GroqCloud, Claude, DeepSeek)} & 09/11/2026 & Completado \\
\hline
3.0 & \multicolumn{4}{p{9.2cm}|}{Anonimizar los datos personales (PII) dentro del dataset de entrenamiento (\textbf{ACT-04}) y verificar cada lote autorizado mediante hashes SHA-256 antes de entrenar al modelo} & 14/11/2026 & Completado \\
\hline
4.0 & \multicolumn{4}{p{9.2cm}|}{Habilitar un asistente de programación privado dentro de la infraestructura física de servidores GPU (\textbf{ACT-05}, 10.0.0.230) para que los desarrolladores trabajen sin enviar código fuera de la red interna} & 22/11/2026 & En progreso \\
\hline
5.0 & \multicolumn{4}{p{9.2cm}|}{Agregar revisión automática de secretos (GitLeaks) y filtros contra \textit{Prompt Injection} en el archivo de Python que es el bot (\textbf{ACT-03}), limitando su cuenta de base de datos (\textbf{ACT-01}) a permisos de solo lectura} & 29/11/2026 & En progreso \\
\hline

% ==================== SECCIÓN V - MONITOREO ====================
\multicolumn{7}{|c|}{\cellcolor{xlHeader}\textbf{Sección V - Monitoreo (RK-002)}} \\
\hline
\multicolumn{2}{|l|}{\textbf{Última evaluación del riesgo}} & \multicolumn{2}{c|}{14/11/2026} & \multicolumn{2}{l|}{\textbf{Siguiente evaluación}} & 14/12/2026 \\
\hline
\multicolumn{7}{|c|}{\textbf{Indicadores clave de riesgo (KRIs)}} \\
\hline
\textbf{\#} & \multicolumn{4}{c|}{\textbf{Indicador}} & \textbf{Medición} & \textbf{Fecha} \\
\hline
1.0 & \multicolumn{4}{p{9.2cm}|}{Intentos mensuales de envío de partes del archivo de Python que es el bot o del dataset hacia sitios externos de IA bloqueados por el filtro DLP y el firewall} & 9 alertas bloqueadas/mes & 14/11/2026 \\
\hline
2.0 & \multicolumn{4}{p{9.2cm}|}{Porcentaje de lotes del dataset de entrenamiento (\textbf{ACT-04}) anonimizados y validados con hash SHA-256} & 100\% de lotes en producción & 14/11/2026 \\
\hline
3.0 & \multicolumn{4}{p{9.2cm}|}{Porcentaje de desarrolladores utilizando el asistente de código local montado en el clúster GPU interno (\textbf{ACT-05})} & 75\% completado & 14/11/2026 \\
\hline
\end{longtable}
\endgroup

% ---------------------------------------------------------
\subsection{Cálculo del Nivel de Riesgo}
% ---------------------------------------------------------

El nivel de riesgo de cada escenario se obtiene multiplicando la frecuencia anual ($F$) por el promedio de los cuatro rubros de impacto ($I_{\text{prom}}$):

\begin{enumerate}[leftmargin=1.5cm]
    \item \textbf{Escenario 1 (RK-001: Intrusión en webserver, escalación a root y pivoteo hacia codeserver):}
    \begin{itemize}[noitemsep,topsep=2pt]
        \item \textbf{Frecuencia Anual ($F_1$):} $F_1 = 5.0$ ($100 < N$ eventos anuales de escaneo y fuerza bruta sobre el servidor expuesto en la DMZ).
        \item \textbf{Calificación General del Impacto ($I_{\text{prom},1}$):}
        \[
            I_{\text{prom},1} = \frac{I_{\text{ganancia}} + I_{\text{costo\_respuesta}} + I_{\text{reputacional}} + I_{\text{multas}}}{4} = \frac{4.0 + 2.0 + 4.0 + 3.0}{4} = \frac{13.0}{4} = 3.25
        \]
        \item \textbf{Calificación General del Nivel de Riesgo ($R_1$):}
        \[
            R_1 = F_1 \times I_{\text{prom},1} = 5.0 \times 3.25 = 16.25 \quad \Longrightarrow \quad \textbf{Nivel de Riesgo: CRÍTICO}
        \]
    \end{itemize}

    \item \textbf{Escenario 2 (RK-002: ShadowAI y exposición del archivo de Python que es el bot, dataset y BD):}
    \begin{itemize}[noitemsep,topsep=2pt]
        \item \textbf{Frecuencia Anual ($F_2$):} $F_2 = 4.0$ ($10 < N \leq 100$ eventos anuales de uso no autorizado de plataformas externas de IA).
        \item \textbf{Calificación General del Impacto ($I_{\text{prom},2}$):}
        \[
            I_{\text{prom},2} = \frac{I_{\text{ganancia}} + I_{\text{costo\_respuesta}} + I_{\text{reputacional}} + I_{\text{multas}}}{4} = \frac{4.0 + 1.0 + 3.0 + 3.0}{4} = \frac{11.0}{4} = 2.75
        \]
        \item \textbf{Calificación General del Nivel de Riesgo ($R_2$):}
        \[
            R_2 = F_2 \times I_{\text{prom},2} = 4.0 \times 2.75 = 11.00 \quad \Longrightarrow \quad \textbf{Nivel de Riesgo: ALTO}
        \]
    \end{itemize}
\end{enumerate}

\newpage
\nocite{*}
\printbibliography[heading=bibintoc, title={Referencias}]

\end{document}
